%%%%%%%%%%%%%%%%%%%%%%%%%%%%%%%%%%%%%%%%%
% Medium Length Professional CV（中文版）
% LaTeX Template (Version 2.0, 8/5/13) — original author: Trey Hunner
% 依赖同目录 resume.cls
% 更新：2026-10-22（华电任职、ICML/ICLR 领域主席、项目金额订正、服务口径统一）
%%%%%%%%%%%%%%%%%%%%%%%%%%%%%%%%%%%%%%%%%

\documentclass{resume}
\hfuzz=100pt

\usepackage[UTF8]{ctex}
\usepackage{xcolor}
\usepackage{fontspec}

\usepackage[backend=biber, style=numeric, sorting=ydnt,
            giveninits=false, maxnames=999, minnames=999]{biblatex}

\AtEveryBibitem{%
  \clearfield{pages}%
  \clearfield{pagination}%
}

% 在每条文献末尾输出 ranking（如 CCF-A）
\renewbibmacro*{finentry}{%
  \iffieldundef{ranking}
    {}
    {\space\textbf{[}\printfield{ranking}\textbf{]}}%
  \finentry
}

\addbibresource{publications.bib}

\definecolor{darkred}{RGB}{139,0,0}
\defbibheading{none}{}
\usepackage[left=0.75in,top=0.6in,right=0.75in,bottom=0.6in]{geometry}
\renewcommand{\refname}{}

\name{刘成}
\address{华北电力大学 控制与计算机工程学院 \\ Email: liucheng@ncepu.edu.cn \quad 电话: +86-18110054096}
\address{地址：北京市昌平区北农路2号，邮编 102206}

\begin{document}

\begin{rSection}{个人简介}
    \vspace{-1em}
\item 刘成，华北电力大学控制与计算机工程学院教授（2026 年 6 月起），此前任中国科学院计算技术研究所处理器芯片全国重点实验室副研究员（2018.06--2026.05）。研究方向为领域专用芯片设计、容错芯片设计与 AI 辅助芯片设计，近期聚焦基于大语言模型的硬件代码生成、代码生成智能体与专用 SoC 自动化设计。主持国家自然科学基金面上/青年项目、中科院先导专项课题、国家重点研发计划子课题、军工专项以及航天科技集团、华为等横向课题，研究成果应用于星载智能计算、状态监测与大数据处理等场景。发表集成电路设计相关学术论文近 100 篇，其中通信/一作 50 余篇，出版英文专著 1 部、书籍章节 2 个，获得\textcolor{darkred}{\textbf{2021 年 IEEE Transactions on Computers（体系结构顶刊，CCF-A）年度唯一最佳论文奖}}以及 GLVLSI 国际会议最佳论文奖。研究成果亦获产业界认可，获 2019 年中科院科技成果转化特等奖、\textcolor{darkred}{\textbf{2023 年华为奥林帕斯先锋奖（国际范围评选，每年 4 项）}}。指导学生获多项国际竞赛前三：ICCAD EDA 竞赛一等奖、DAC 系统设计竞赛二/三等奖、CVPR 低功耗计算机视觉挑战赛三等奖、PPoPP 2025 图计算竞赛三等奖。曾多次获中科院计算所优秀员工称号。

担任 IEEE Transactions on Emerging Topics in Computing (TETC) 编委、IEEE Transactions on Parallel and Distributed Systems (TPDS) 评审委员；任 \textbf{ICML 2027、ICLR 2027 领域主席（Area Chair）}；担任 DAC、FPT、ATS、DFT、DFTS 等设计自动化与 FPGA 领域国际会议的程序委员，并担任 NeurIPS、ICML、ICLR、AAAI 等 AI 会议审稿人及 TCAD、TVLSI、TPDS 等领域期刊审稿人。现为 IEEE 高级会员、CCF 高级会员。
    \vspace{-0.5em}
\end{rSection}

\begin{rSection}{教育与工作经历}
{\bf 华北电力大学} \hfill {\em 2026-06 - 至今} \\
控制与计算机工程学院，教授\\[0.3em]
{\bf 中国科学院计算技术研究所} \hfill {\em 2018-06 - 2026-05} \\
处理器芯片全国重点实验室（原计算机体系结构国家重点实验室），副研究员\\[0.3em]
{\bf 新加坡国立大学} \hfill {\em 2016-12 - 2018-06} \\
计算机学院，博士后\\[0.3em]
{\bf 香港大学} \hfill {\em 2011-09 - 2016-11} \\
计算机工程专业，电机与电子工程系，博士\\[0.3em]
{\bf 哈尔滨工业大学} \hfill {\em 2007-09 - 2009-07} \\
微电子学与固体电子学专业，航天学院，硕士\\[0.3em]
{\bf 哈尔滨工业大学} \hfill {\em 2003-09 - 2007-07} \\
电子信息科学与技术专业，航天学院，学士
\end{rSection}

\begin{rSection}{科研项目}
{\bf 近似计算范式下的领域专用芯片容错设计方法研究} \hfill {\em 2027 - 2030} \\
国家自然科学基金面上项目，60.6 万，主持（2026.8 获批）\\[0.3em]
{\bf 面向处理器芯片跨层优化的软硬件平台} \hfill {\em 2024 - 2028} \\
中科院先导 B 计划课题，1074 万，主持\\[0.3em]
{\bf 面向集成电路多模态数据的表示与应用研究} \hfill {\em 2025 - 2026} \\
中科院算力专项，71 万，主持\\[0.3em]
{\bf 弹性容错深度学习处理器设计} \hfill {\em 2022 - 2025} \\
国家自然科学基金面上项目，57 万，主持\\[0.3em]
{\bf 基于 FPGA 的高能效图计算加速器设计} \hfill {\em 2020 - 2022} \\
国家自然科学基金青年项目，27 万，主持\\[0.3em]
{\bf 智能微系统可靠性测评方法研究} \hfill {\em 2023 - 2024} \\
航天科技集团横向，50 万，主持\\[0.3em]
{\bf 容错深度学习工具链研究} \hfill {\em 2022 - 2023} \\
航天科技集团横向，30 万，主持\\[0.3em]
{\bf 智能存算融合架构和任务调度算法研究} \hfill {\em 2023 - 2024} \\
军工专项，130 万，主持\\[0.3em]
{\bf 领域专用加速系统与示范应用} \hfill {\em 2022 - 2025} \\
国家重点研发计划子课题，160 万，子课题负责人\\[0.3em]
{\bf 基于智能网卡与智能存储设备的流式计算系统研究} \hfill {\em 2020 - 2021} \\
华为横向，120 万，主持\\[0.3em]
{\bf 深度学习处理器自动化容错设计方法} \hfill {\em 2021 - 2022} \\
计算机体系结构重点实验室培育项目，30 万，主持\\[0.3em]
{\bf 人工智能芯片测评技术研究} \hfill {\em 2023 - 2024} \\
专项项目，290 万，参与\\[0.3em]
{\bf 基于分离式内存的近存智能计算架构研究} \hfill {\em 2024 - 2025} \\
专项项目，130 万，参与\\[0.3em]
{\bf 基于微存储的计算型存储} \hfill {\em 2023 - 2024} \\
华为横向，120 万，参与\\[0.3em]
{\bf 开源物端智能处理器} \hfill {\em 2020 - 2021} \\
中科院先导 C 课题，1200 万，参与
\end{rSection}

\begin{rSection}{论文}
\end{rSection}
\nocite{*}
\printbibliography[heading=none]

\end{document}
